%! Populations, Samples, Parameters, and Statistics — Unit 1, Lesson 1.2 — Warm-Up (Answer Key)
% ONE PAGE, blank and key. Three items at 12pt (the stats-model size; replaces 1.1's \large-at-10pt trick),
% \workrowsep 50pt, item spacing 22pt) so each item gets real handwriting room and still fits.
% Do not add a fourth item — the page is full.
% GRADUAL RELEASE: the warm-up activates prior knowledge before the guided notes. Items 1 and 2
% are last night's Lakeside arithmetic — the same 50-of-800 survey the notes open with, so the
% survey students meet in section 1 is already half-computed. Item 3 makes them separate the
% two groups WITHOUT the vocabulary: section 1 names them POPULATION and SAMPLE ten minutes
% later, and the second half of section 1 names the 52% a STATISTIC.
\documentclass[12pt]{article}
\usepackage{saar-article}
\usepackage{saar-key}   % pulls in -boxes; do NOT also load -boxes
\newcommand{\TallMath}[1]{$\displaystyle #1\rule[-1.4em]{0pt}{3.2em}$}
% Handwriting room inside every \begin{work} block. Moves the blank and the
% key together, so the two can never drift apart.
\setlength{\workrowsep}{50pt}

\begin{document}

\pageheader{Unit 1, Lesson 1.2}{Warm-Up --- Answer Key}
% NAMESTRIP: no \namedateperiod here — mirrors the blank. NO teachernote here —
% teacher prose lives in the lesson plan.

\vspace{0.15in}

\begin{notesbox}{Spiral Review --- Back to the Lakeside Farmers Market}
\normalsize
About $800$ shoppers come to the Lakeside Farmers Market on a Saturday morning.
One Saturday the manager asked $50$ of them how they traveled to the market.
Of the $50$ she asked, $26$ drove.

\begin{enumerate}[leftmargin=1.8em, itemsep=22pt, topsep=10pt]

  \item What \textbf{percent} of the shoppers she asked drove to the market?

        \begin{work}
          \dfrac{26}{50} &= 0.52 \\
                         &= 52\%
        \end{work}

  \item Based on that survey, \textbf{predict} how many of the $800$ Saturday
        shoppers drove.

        \begin{work}
          0.52 \times 800 &= 416 \text{ shoppers}
        \end{work}

  \item Two different groups of shoppers appear in the paragraph above. Write
        how many are in each group.

        \vspace{4pt}
        \renewcommand{\arraystretch}{2.6}
        \begin{tabularx}{\linewidth}{@{} X p{3.4cm} @{}}
        \toprule
        \textbf{Group} & \textbf{How many} \\
        \midrule
        The group the manager actually collected data from & \ans{50} \\
        The whole group her report is about                & \ans{800} \\
        \bottomrule
        \end{tabularx}

        \vspace{10pt}
        Which number did she \emph{count} from real answers --- the $52\%$ or
        the $416$? \ans{the 52\%}

\end{enumerate}
\end{notesbox}

\end{document}
